\documentclass[14pt]{extarticle}
\usepackage{fontspec}
\usepackage{polyglossia}
\usepackage{amsmath}
\usepackage{amssymb}
\usepackage{xcolor}
\usepackage{fancyhdr}
\usepackage{graphicx}
\usepackage{listings}
\usepackage[most]{tcolorbox}
% \usepackage{setspace}
\usepackage{pifont}
\usepackage{enumitem}
\usepackage{titlesec}
\usepackage[bottom]{footmisc}
\usepackage{titling}
\usepackage{minted}
\usepackage{etoolbox}
\usepackage{array}
\usepackage{longtable} % page-breakable tables (must load before polyglossia/bidi RTL setup)
\usepackage{extsizes}
\usepackage{float}
\usepackage{booktabs}
\usepackage{hyperref}
\usepackage[onehalfspacing]{setspace}
\newif\ifwhiteblack
\ifdefined\setwhiteblack
  \setwhiteblack
\fi

\ifwhiteblack
  \usepackage[a4paper, landscape, top=1.25cm, bottom=1.25cm, left=1.25cm, right=1.25cm, headheight=1.5cm, headsep=0.5cm]{geometry}
\else
  \usepackage[a4paper, top=2.7cm, bottom=2.7cm, left=2.5cm, right=2.5cm, headheight=1.5cm, headsep=0.5cm]{geometry}
\fi
\usepackage{tikz}
\usepackage{datetime2}

\usetikzlibrary{automata, positioning, arrows.meta, arrows}
\providecommand{\LTcaptype}{table}

\newfontfamily\emoji{Segoe UI Emoji}

\pagestyle{fancy}

\setmainlanguage[numerals=western]{arabic}
\setotherlanguage{english}
\newfontfamily\arabicfont[Script=Arabic]{Calibri}
\newfontfamily\arabicfonttt[Script=Arabic]{Courier New}



\ifwhiteblack
\lstset{
  language=[Sharp]C,
  numbers=left,
  stepnumber=1,
  numbersep=8pt,
  frame=single,
  basicstyle=\ttfamily\small,
  keywordstyle=\bfseries,
  stringstyle=\ttfamily,
  commentstyle=\itshape
}
\else
\lstset{
  language=[Sharp]C,
  numbers=left,
  stepnumber=1,
  numbersep=8pt,
  frame=single,
  basicstyle=\ttfamily\small,
  keywordstyle=\color{blue},
  stringstyle=\color{red},
  commentstyle=\color{green!50!black}
}
\fi

\newif\ifdetailed
\ifdefined\setdetailed
  \setdetailed
\fi

\newif\ifwithsols
\ifdefined\setwithsols
  \setwithsols
\fi

\ifwhiteblack
  \AtBeginDocument{
    \LARGE
    \usemintedstyle{bw}
  }
\fi

% Default minted look: compact code that wraps long lines.
% Not applied in black-and-white landscape mode, which already uses \LARGE.
\ifwhiteblack\else
  \setminted{fontsize=\small, breaklines=true}
\fi

% unified tcolorboxes for articles
\ifwhiteblack
\tcbset{colback=white, colframe=black, fonttitle=\bfseries, boxrule=1pt}
\newtcolorbox{boxDef}[1][]{title={{\emoji📘} تعريف\ifx\\#1\\\else ~#1\fi :}}
\newtcolorbox{boxExercise}[1][]{title={{\emoji🧩} تمرين\ifx\\#1\\\else ~#1\fi :}}
\newtcolorbox{boxExample}[1][]{title={{\emoji📝} مثال\ifx\\#1\\\else ~#1\fi :}}
\newtcolorbox{boxNote}[1][]{title={{\emoji✨} ملاحظة\ifx\\#1\\\else ~#1\fi :}}
\newtcolorbox{boxAttention}[1][]{title={{\emoji🔔} تنبيه\ifx\\#1\\\else ~#1\fi :}}
\newtcolorbox{boxWarning}[1][]{title={{\emoji⚡} ملاحظة هامة\ifx\\#1\\\else ~#1\fi :}}
\newtcolorbox{boxSolution}[1][]{title={{\emoji✅} حل\ifx\\#1\\\else ~#1\fi :}}
\newtcolorbox{boxSymbol}[1][]{title={{\emoji🔣} رمز\ifx\\#1\\\else ~#1\fi :}}
\newtcolorbox{boxHint}[1][]{title={{\emoji💡} تلميح\ifx\\#1\\\else ~#1\fi :}}
\else
\tcbset{colback=white, colframe=black, fonttitle=\bfseries, boxrule=0.8pt}
\newtcolorbox{boxDef}[1][]{colback=blue!5!white,colframe=blue!75!black,
  title={{\emoji📘} تعريف\ifx\\#1\\\else ~#1\fi :}}
\newtcolorbox{boxExercise}[1][]{colback=cyan!5!white,colframe=cyan!70!black,
  title={{\emoji🧩} تمرين\ifx\\#1\\\else ~#1\fi :}}
\newtcolorbox{boxExample}[1][]{colback=yellow!5!white,colframe=orange!90!black,
  title={{\emoji📝} مثال\ifx\\#1\\\else ~#1\fi :}}
\newtcolorbox{boxNote}[1][]{colback=gray!10!white,colframe=black,
  title={{\emoji✨} ملاحظة\ifx\\#1\\\else ~#1\fi :}}
\newtcolorbox{boxAttention}[1][]{colback=magenta!10!white,colframe=magenta!80!black,
  title={{\emoji🔔} تنبيه\ifx\\#1\\\else ~#1\fi :}}
\newtcolorbox{boxWarning}[1][]{colback=red!5!white,colframe=red!75!black,
  title={{\emoji⚡} ملاحظة هامة\ifx\\#1\\\else ~#1\fi :}}
\newtcolorbox{boxSolution}[1][]{colback=green!5!white,colframe=green!60!black,
  title={{\emoji✅} حل\ifx\\#1\\\else ~#1\fi :}}
\newtcolorbox{boxSymbol}[1][]{colback=purple!5!white,colframe=purple!70!black,
  title={{\emoji🔣} رمز\ifx\\#1\\\else ~#1\fi :}}
\newtcolorbox{boxHint}[1][]{colback=teal!5!white,colframe=teal!60!black,
  title={{\emoji💡} تلميح\ifx\\#1\\\else ~#1\fi :}}
\fi


\ifwhiteblack
\tcbset{simplecode/.style={ colback=white, colframe=black, boxrule=1pt, arc=2pt, left=4pt,right=4pt,top=4pt,bottom=4pt}}
\else
\tcbset{simplecode/.style={ colback=gray!5, colframe=black!50, boxrule=0.4pt, arc=2pt, left=4pt,right=4pt,top=4pt,bottom=4pt}}
\fi
% Compact code boxes: tighter line spacing (the preamble uses onehalfspacing) and small vertical gaps.
% Must come AFTER the simplecode definition above.
\tcbset{simplecode/.append style={fontupper=\setstretch{0.95}, top=4pt, bottom=4pt, before skip=6pt, after skip=6pt}}
\newenvironment{boxCode}{\begin{tcolorbox}[simplecode]}{\end{tcolorbox}}

% Code cell used inside tables (keeps code left-to-right inside RTL tables).
% Use inside a fixed-width p{...} column: \begin{pycell}\begin{minted}...\end{minted}\end{pycell}
\newtcolorbox{pycell}{frame empty,colback=white,boxsep=0pt,left=2pt,right=2pt,top=2pt,bottom=2pt,width=\linewidth}

\newcolumntype{C}[1]{>{\centering\arraybackslash}p{#1}}

\makeatletter
\DTMnewdatestyle{dotted}{%
  \renewcommand{\DTMdisplaydate}[4]{\two@digits{##3}.\two@digits{##2}.##1}%
  \renewcommand{\DTMDisplaydate}{\DTMdisplaydate}%
}
\makeatother
\DTMsetdatestyle{dotted}

% redefine spaces after titles
\makeatletter
\renewcommand{\@maketitle}{%
  \begin{center}
    {\huge \bfseries \@title \par}%
    \vskip 0.2em % space between title and author
    {\large \@author \par}%
    % \vskip 0.2em % space between author and date
    % {\normalsize \@date \par}%
    {\small \hfill \begin{english}\DTMtoday\end{english} \par}%
  \end{center}
}
\makeatother

\fancyhf{} % clear default
\fancypagestyle{plain}{
  \fancyhf{}
  \fancyhead[L]{مدرسة التسامح الشاملة}
  % \fancyhead[L]{\includegraphics[height=1cm]{../../../images/logoTasamoh.png}}
  \fancyhead[R]{الأستاذ محمود اغبارية}
  \fancyfoot[C]{\thepage}
}

\fancyhead[L]{مدرسة التسامح الشاملة}
\fancyhead[R]{الأستاذ محمود اغبارية}
\fancyfoot[C]{\thepage}

\setcounter{tocdepth}{3} % only section subsection and subsubsection in TOC

\newcommand{\importpdfpage}[2]{%
    \thispagestyle{empty}
    \newgeometry{left=0cm,right=0cm,top=0cm,bottom=0cm}
    \includegraphics[page=#2,width=0.95\paperwidth,height=0.95\paperheight,keepaspectratio]{#1}
    \restoregeometry
}

\newcommand{\insertFullImg}[1]{%
    \noindent
    \makebox[\textwidth][c]{
        \includegraphics[width=0.9\paperwidth,keepaspectratio]{#1}%
    }%
}%



\ifwithsols
\title{حل وظيفة بيتية 5 للصف العاشر 10: $\mathtt{if\ Statement}$}
\else
\title{وظيفة بيتية 5 للصف العاشر 10: $\mathtt{if\ Statement}$}
\fi


\begin{document}

\maketitle
\thispagestyle{fancy}

\begin{enumerate}[itemsep=3em]

\item
اكتب برنامجًا يستقبل من المستخدم عددًا صحيحًا ثنائيّ المنازل (بين 10 و 99)، ويطبع "\textbf{\textenglish{Yes}}" إذا كانت منزلتاه عددين متتاليين (أي أنّ الفرق المطلق بينهما يساوي 1، مثل 34 أو 43)، وإلا يطبع "\textbf{\textenglish{No}}".
\ifdetailed
\begin{boxExample}[1]
\begin{english}
\begin{minted}{text}
Enter a two-digit number:
34
Yes
\end{minted}
\end{english}
\end{boxExample}
\begin{boxExample}[2]
\begin{english}
\begin{minted}{text}
Enter a two-digit number:
79
No
\end{minted}
\end{english}
\end{boxExample}
\fi
\ifwithsols
\begin{boxSolution}
\begin{english}
\begin{minted}{csharp}
Console.WriteLine("Enter a two-digit number:");
int num = int.Parse(Console.ReadLine());
int tens = num / 10;
int units = num % 10;
if (Math.Abs(tens - units) == 1)
{
    Console.WriteLine("Yes");
}
else
{
    Console.WriteLine("No");
}
\end{minted}
\end{english}
\end{boxSolution}
\fi
\clearpage

\item
اكتب برنامجًا يستقبل من المستخدم عددًا صحيحًا ثلاثيّ المنازل (بين 100 و 999)، ويطبع "\textbf{\textenglish{Palindrom}}" إذا كانت قراءته من اليمين إلى اليسار متطابقة مع قراءته من اليسار إلى اليمين (مثلاً: 121، 585).
\ifdetailed
\begin{boxExample}[1]
\begin{english}
\begin{minted}{text}
Enter a three-digit number:
121
Palindrom
\end{minted}
\end{english}
\end{boxExample}
\begin{boxExample}[2]
\begin{english}
\begin{minted}{text}
Enter a three-digit number:
123
\end{minted}
\end{english}
\end{boxExample}
\fi
\ifwithsols
\begin{boxSolution}
\begin{english}
\begin{minted}{csharp}
Console.WriteLine("Enter a three-digit number:");
int num = int.Parse(Console.ReadLine());
int hundreds = num / 100;
int units = num % 10;
if (hundreds == units)
{
    Console.WriteLine("Palindrom");
}
\end{minted}
\end{english}
\end{boxSolution}
\fi
\clearpage

\item
اكتب برنامجًا يستقبل من المستخدم 3 أعداد عشرية بترتيب غير تنازلي (كل عدد أكبر أو يساوي الذي قبله: $a \le b \le c$) تمثل أطوال أضلاع مثلث. \\
على البرنامج أن يطبع نوع المثلث:
\begin{itemize}
    \item "\textbf{\textenglish{Equilateral}}" (مثلث متساوي الأضلاع)، إذا كانت كل أضلاعه متساوية.
    \item "\textbf{\textenglish{Isosceles}}" (مثلث متساوي الساقين)، إذا كان هناك ضلعان بالضبط متساويان.
    \item "\textbf{\textenglish{Right triangle}}" (مثلث قائم الزاوية)، إذا كان المثلث يحقق قانون فيثاغورس: $a^2 + b^2 = c^2$ ($c$ هو الضلع الأطول بين الثلاثة).
\end{itemize}
\ifdetailed
\begin{boxExample}[1]
\begin{english}
\begin{minted}{text}
Enter three sides:
3
3
3
Equilateral
\end{minted}
\end{english}
\end{boxExample}
\begin{boxExample}[2]
\begin{english}
\begin{minted}{text}
Enter three sides:
3
3
5
Isosceles
\end{minted}
\end{english}
\end{boxExample}
\begin{boxExample}[3]
\begin{english}
\begin{minted}{text}
Enter three sides:
3
4
5
Right triangle
\end{minted}
\end{english}
\end{boxExample}
\fi
\ifwithsols
\begin{boxSolution}
\begin{english}
\begin{minted}{csharp}
Console.WriteLine("Enter three sides:");
double a = double.Parse(Console.ReadLine());
double b = double.Parse(Console.ReadLine());
double c = double.Parse(Console.ReadLine());

// Since a <= b <= c, if a == c then all three sides are equal (a == b == c)
if (a == c)
{
    Console.WriteLine("Equilateral");
}
else if (a == b)
{
    Console.WriteLine("Isosceles");
}
else if (b == c)
{
    Console.WriteLine("Isosceles");
}
else if (Math.Pow(a, 2) + Math.Pow(b, 2) == Math.Pow(c, 2))
{
    Console.WriteLine("Right triangle");
}
\end{minted}
\end{english}
\end{boxSolution}
\fi
\clearpage

% \item
% اكتب برنامجًا يمثّل آلة حاسبة بسيطة: يستقبل البرنامج من المستخدم عددًا، ثم حرفًا (\textenglish{char}) يمثّل العملية الحسابية (\textenglish{'+'}, \textenglish{'-'}, \textenglish{'*'}, \textenglish{'/'}, \textenglish{'\%'})، ثم عددًا آخر، ويحسب ويطبع نتيجة العملية الحسابية.
% \ifdetailed
% \begin{boxExample}[1]
% \begin{english}
% \begin{minted}{text}
% Enter first number:
% 10
% Enter operation (+, -, *, /, %):
% +
% Enter second number:
% 5
% Result: 15
% \end{minted}
% \end{english}
% \end{boxExample}
% \begin{boxExample}[2]
% \begin{english}
% \begin{minted}{text}
% Enter first number:
% 17
% Enter operation (+, -, *, /, %):
% %
% Enter second number:
% 5
% Result: 2
% \end{minted}
% \end{english}
% \end{boxExample}
% \fi

% \ifwithsols
% \begin{boxSolution}
% \begin{english}
% \begin{minted}{csharp}
% Console.WriteLine("Enter first number:");
% double num1 = double.Parse(Console.ReadLine());
% Console.WriteLine("Enter operation (+, -, *, /, %):");
% char op = char.Parse(Console.ReadLine());
% Console.WriteLine("Enter second number:");
% double num2 = double.Parse(Console.ReadLine());

% if (op == '+')
% {
%     Console.WriteLine("Result: " + (num1 + num2));
% }
% else if (op == '-')
% {
%     Console.WriteLine("Result: " + (num1 - num2));
% }
% else if (op == '*')
% {
%     Console.WriteLine("Result: " + (num1 * num2));
% }
% else if (op == '/')
% {
%     Console.WriteLine("Result: " + (num1 / num2));
% }
% else if (op == '%')
% {
%     Console.WriteLine("Result: " + (num1 % num2));
% }
% \end{minted}
% \end{english}
% \end{boxSolution}
% \clearpage
% \fi

\end{enumerate}


\vspace{1cm}
\begin{flushleft}
أرجو لكم وقتًا ممتعًا.

الأستاذ محمود اغبارية.
\end{flushleft}


\end{document}